\begin{MBSourceBox}
\refstepcounter{booksummary}\label{box:homological-algebra-summary}
\noindent\textbf{Box~\thebooksummary:} {Chapter~\ref{chap:homological-algebra} 同调代数内容摘要.}\par

\emph{链复形}是如下一列向量空间与算子:
\[
\cdots\to V_{k+1}\xrightarrow{\partial_{k+1}}V_k\xrightarrow{\partial_k}V_{k-1}\to\cdots
\]

其中相邻算子的复合为零. 它在每个层次 $k$ 上产生三个向量空间:
\begin{itemize}
\item {$\mathcal Z_k$, 即边界算子 $\partial_k$ 的零空间, 由\emph{闭链}组成;}
\item {$\mathcal B_k$, 即 $\partial_{k+1}$ 的像空间, 由\emph{边界链}组成;}
\item {$\mathcal H_k=\mathcal Z_k/\mathcal B_k$, 即\emph{同调空间}.}
\end{itemize}

一个重要例子是与剖分或更一般的单纯复形相伴的\emph{单纯链复形}. 它的同调空间的维数是底空间的 Betti 数.

\emph{上链复形}的指标次序相反,
\[
\cdots\to V^{k-1}\xrightarrow{d^{k-1}}V^k\xrightarrow{d^k}V^{k+1}\to\cdots,
\]

由此产生\emph{闭上链}, \emph{上边界链}与\emph{上同调}.

链复形的对偶空间和对偶映射构成一个上链复形.

另一个上链复形是 \emph{de Rham 复形}, 在三维区域上, 它可以写成
\[
0\to C^\infty(\Omega)\xrightarrow{\operatorname{grad}}C^\infty(\Omega;\mathbb R^3)\xrightarrow{\operatorname{curl}}C^\infty(\Omega;\mathbb R^3)\xrightarrow{\operatorname{div}}C^\infty(\Omega)\to0.
\]

de Rham 映射是从 de Rham 复形到单纯链复形的对偶复形的一个\emph{上链映射}, 因而在上同调上诱导一个映射. De Rham 定理断言, 这一诱导映射是同构.
\end{MBSourceBox}
